Ահա նկարներից դուրս բերված հայերեն տեքստը և մաթեմատիկական բանաձևերը։

Բերենք նաև մաթեմատիկական անալիզի դասընթացից հայտնի ֆունկցիայի հավասարաչափ անընդհատության մասին Կանտորի թեորեմի ընդհանրացումը մետրիկական տարածությունների դեպքում։

<center>166</center>

**Սահմանում։** Դիցուք $X$-ը և $Y$-ը մետրիկական տարածություններ են համապատասխանաբար $\rho_1$ և $\rho_2$ մետրիկաներով, և ունենք $f: X \to Y$ արտապատկերում։ Ապա $f$-ը կոչվում է <u>հավասարաչափ անընդհատ արտապատկերում</u>, եթե ցանկացած $\varepsilon > 0$ թվի դեպքում գոյություն ունի $\delta > 0$ թիվ, որ ամեն մի $x_1, x_2 \in X$ կետերի դեպքում եթե $\rho_1(x_1, x_2) < \delta$, ապա $\rho_2(f(x_1), f(x_2)) < \varepsilon$։

Պարզ է, որ ամեն մի հավասարաչափ անընդհատ արտապատկերում անընդհատ արտապատկերում է։ Հակառակը ընդհանուր դեպքում ճիշտ չէ (<u>հիմնավորեք օրինակով</u>)։

**Թեորեմ 10։** Եթե $(X, \rho_1), (Y, \rho_2)$ մետրիկական տարածություններից առաջինը կոմպակտ տարածություն է, ապա ամեն մի $f: X \to Y$ անընդհատ արտապատկերում հավասարաչափ անընդհատ է։

**Ապացուցում։** Դիցուք $\varepsilon > 0$ թիվը ընտրված է։ Ամեն մի $x \in X$ կետի համար գոյություն ունի $\delta(x) > 0$ թիվ, որ եթե $x' \in X$ և $\rho_1(x', x) < 2\delta(x)$, ապա $\rho_2(f(x'), f(x)) < \varepsilon/2$։ Պարզ է, որ $x$ կենտրոնով և $\delta(x)$ շառավղով $B(x; \delta(x)) = \{x' ; \rho_1(x', x) < \delta(x)\}$ գնդերի $\{B(x, \delta(x)); x \in X\}$ ընտանիքը կազմում է $X$ տարածության բաց ծածկույթ։ Քանի որ տարածությունը կոմպակտ է, գոյություն ունի այդ ծածկույթի $B(x_1, \delta(x_1)), B(x_2, \delta(x_2)), \dots, B(x_n, \delta(x_n))$ վերջավոր ենթածածկույթ։ Նշանակենք $\delta = \min\{\delta(x_1), \delta(x_2), \dots, \delta(x_n)\}$։ Պարզ է, որ եթե $x, y \in X$ և $\rho_1(x, y) < \delta$, ապա $x \in B(x_i, \delta(x_i))$ ինչ-որ $i (1 \le i \le n)$ ինդեքսի դեպքում։ Այնուհետև, քանի որ


$$\rho_1(y, x_i) \le \rho_1(y, x) + \rho_1(x, x_i) < \delta + \delta(x_i) \le 2\delta(x_i)$$


ուստի $\rho_2(f(y), f(x_i)) < \varepsilon/2$։ Այստեղից ստանում ենք՝


$$\rho_2(f(x), f(y)) \le \rho_2(f(x), f(x_i)) + \rho_2(f(x_i), f(y)) < \varepsilon: \blacksquare$$


Ահա նկարից դուրս բերված հայերեն տեքստը.

...ընդհանուր տոպոլոգիայում և նրա կիրառություններում դիտարկվում են կոմպակտության բնույթի, կամ կոմպակտության հետ առնչվող մի շարք հասկացություններ ինչպիսիք են հաշվելի-կոմպակտ, սեկվենցիալ կոմպակտ, լոկալ կոմպակտ, պարակոմպակտ տարածություն հասկացությունները։ Ցանկանում ենք սրանց շարքում նշել նաև թեմա 8-ից մեզ ծանոթ Լինդելյոֆյան տարածություն հասկացությունը։

**Սահմանում։** Տարածությունը կոչվում է <u>հաշվելի-կոմպակտ տարածություն</u>, եթե նրա ցանկացած հաշվելի բաց ծածկույթից կարելի է անջատել

<center>167</center>

վերջավոր ենթածածկույթ։


Ահա վերջին երկու նկարներից դուրս բերված տեքստը և բանաձևերը.

Պարզ է, որ ամեն մի կոմպակտ տարածություն նաև հաշվելի-կոմպակտ տարածություն է։ Հակառակ պնդումը ընդհանուր դեպքում ճիշտ չէ (կան տոպոլոգիական տարածությունների մի քանի օրինակներ, երբ հաշվելի-կոմպակտ տարածությունը կոմպակտ տարածություն չէ)։ Այդ բոլոր օրինակները բավականին բարդ են, քանի որ կոմպակտության այս երկու հասկացությունների միջև տարբերությունը այնքան էլ մեծ չէ, ինչը անուղղակի ձևով հաստատվում է հետևյալ պնդումով։

**Թեորեմ 11։** Հաշվելիության երկրորդ աքսիոմին բավարարող տարածությունների դեպքում կոմպակտությունը համարժեք է հաշվելի-կոմպակտությանը։

**Ապացուցում։** Դիցուք $X$-ը հաշվելի-կոմպակտ տարածություն է, և $S = \{U_\alpha; \alpha \in A\}$ ընտանիքը նրա կամայական բաց ծածկույթ է։ Թեորեմի պայմանից ունենք, որ $X$-ը Լինդելյոֆյան տարածություն է (ըստ սահմանման)։ Այժմ թեմա 8-ի թեորեմ 3-ից հետևում է՝ $S$ ծածկույթից կարելի է անջատել $X$-ի հաշվելի $S'$ ենթածածկույթ, որն էլ իր հերթին պարունակում է $S''$ վերջավոր ենթածածկույթ։ Պարզ է, որ $S''$-ը վերջավոր ենթածածկույթ է նաև $S$ ծածկույթի համար, ինչով և ավարտվում է թեորեմի ապացուցումը։

---

**Սահմանում։** Տոպոլոգիական $X$ տարածությունը կոչվում է <u>սեկվենցիալ-կոմպակտ տարածություն</u>, եթե այն բավարարում է Բոլցանո-Վեյերշտրասի պայմանին՝ $X$-ում կետերի ցանկացած անվերջ հաջորդականություն պարունակում է զուգամետ ենթահաջորդականություն (նշենք՝ սեկվենցիալ բառը լատիներենից թարգմանաբար նշանակում է հաջորդականություն)։

Սեկվենցիալ-կոմպակտ տարածության պարզ օրինակ է թվային ուղղի կամայական $[a, b]$ հատվածը։ Մաթեմատիկական անալիզի դասընթացում հատվածի կիսման եղանակով ցույց է տրվում որ $[a, b]$-ն բավարարում է սահմանման պահանջին։

**Թեորեմ 12։** Մետրիկական տարածությունների, մասնավորապես $\mathbb{R}^n$ Էվկլիդեսյան տարածությունների դեպքում կոմպակտությունը, հաշվելի-կոմպակտությունը և սեկվենցիալ կոմպակտությունը համարժեք են միմյանց։


Ահա վերջին երկու նկարներից (`image_ed0794.jpg` և `image_ed0775.jpg`) դուրս բերված տեքստը և բանաձևերը։

...Ապացույցի, ինչպես նաև կոմպակտության հետ առնչվող այլ հասկացությունների հետ կարելի է ծանոթանալ [ ] գրքում։ Իսկ մենք այստեղ սահմանափակվենք նշելով մետրիկական...

**Օրինակ 3։** Ամեն մի $(X, \text{վերջ. co})$ տարածություն կոմպակտ տարածություն է։ Իրոք, դիցուք $S = \{U_i; i \in I\}$ ընտանիքը $X$-ի որևէ բաց ծածկույթ է, ընդ որում $U_i$ ենթաբազմություններից ոչ մեկը չի համընկնում $X$-ի հետ (հակառակ դեպքում ապացուցելու բան չկա)։ Դիտարկենք այդ ծածկույթի որևէ $U_{i_0}$ տարր և նրա $X \setminus U_{i_0} = \{x_1, x_2, \dots, x_n\}$ լրացումը $X$-ում։ Այժմ յուրաքանչյուր $x_k, k=1, 2, \dots, n$ տարրի համար ընտրելով որևէ $U_{i_k} \in S$ տարր այնպես, որ $x_k \in U_{i_k}$, կստանանք $X$-ի $S$ ծածկույթի $U_{i_0}, U_{i_1}, \dots, U_{i_n}$ վերջավոր ենթածածկույթը։ (Տես նաև էջ 162)

...բազմությունների կոմպակտ ենթաբազմությունների ևս մի քանի հասկացություններ։

**Սահմանում։** Դիցուք $A$-ն և $M$-ը $(X, \rho)$ մետրիկական տարածության ենթաբազմություններ են, ընդ որում $A \subset M$։ Ասում են, որ $A$-ն $\varepsilon$-խիտ է $M$-ում, եթե $M$-ի ցանկացած $x$ կետի համար գոյություն ունի $a \in A$ կետ, որ $\rho(x, a) < \varepsilon$։

---

Ցույց տանք, որ $(X, \rho)$-ի ամեն մի $M$ կոմպակտ ենթաբազմության համար գոյություն ունի նրա $\varepsilon$-խիտ վերջավոր ենթաբազմություն, որտեղ $\varepsilon$-ը կամայական սևեռված դրական թիվ է։

Դիտարկենք $M$-ի $\{B(x, \varepsilon); x \in M\}$ բաց ծածկույթը կազմված $x$ կենտրոններով և $\varepsilon$ շառավղով $B(x, \varepsilon)$ գնդերից։ Օգտվելով $M$-ի կոմպակտությունից ընտրենք այդ ծածկույթի որևէ $S'(\varepsilon) = \{B(a_1, \varepsilon), B(a_2, \varepsilon), \dots, B(a_k, \varepsilon)\}$ վերջավոր ենթածածկույթ։ Պարզ է, որ $M$-ի $A(\varepsilon) = \{a_1, a_2, \dots, a_k\}$ ենթաբազմությունը որոնելին է (հիմնավորեք)։

**Թեորեմ 13։** Մետրիկական տարածության ամեն մի $M$ կոմպակտ ենթաբազմություն սեպարաբել է, ուստի և բավարարում է հաշվելիության երկրորդ աքսիոմին։

**Ապացուցում։** Ամեն մի $n$ բնական թվի համար ընտրելով $M$-ի որևէ $A(1/n)$ վերջավոր $1/n$-խիտ ենթաբազմություն և դիտարկենք $A = \bigcup A(1/n)$ միավորումը։ Ստուգումը, որ $A$-ն $M$-ի հաշվելի ամենուրեք խիտ ենթաբազմություն է թողնված է ընթերցողին։ $\blacksquare$

<center>169</center>